从水平的边际点图可以清楚看出, 这两组数据在 $x_1$ 值中有较多的重叠. 类似的, 垂直的边际点图显示出两组数据在 $x_2$ 值中存在较多重叠. 这些散布图暗示, 对每一个组来说, 两个变量之间都有相当强的正相关, 而且虽然有重叠现象发生, 但组 1 的测量数据大多位于组 2 的数据的东南方.

设 $\boldsymbol{\mu}_1'=[\mu_{11},\mu_{12}]$ 为第一组的总体均值向量, $\boldsymbol{\mu}_2'=[\mu_{21},\mu_{22}]$ 为第二组的总体均值向量, 利用 $x_1$ 的观测值, 由一元方差分析知 $F=2.46$, 且其自由度为 $\nu_1=1$ 和 $\nu_2=18$. 故在任何合理的显著性水平下 $[F_{1,18}(0.10)=3.01]$, 我们不能拒绝假设 $H_0:\mu_{11}=\mu_{21}$. 而利用 $x_2$ 的观测值, 由一元方差分析知 $F=2.68$, 且其自由度为 $\nu_1=1$ 和 $\nu_2=18$, 同样在任何合理的显著性水平下, 我们不能拒绝假设 $H_0:\mu_{12}=\mu_{22}$. 故一元检验表明两组总体的均值分量间无显著差异, 因而得出 $\boldsymbol{\mu}_1=\boldsymbol{\mu}_2$ 的结论.

另一方面, 如果我们用霍特林 $T^2$ 来检验均值向量的相等性, 则可以发现 $T^2=17.29>c^2=\frac{18\times 2}{7}F_{2,17}(0.01)=2.118\times 6.11=12.94$, 因此, 在显著性水平 $1\%$ 下, 我们拒绝 $H_0:\boldsymbol{\mu}_1=\boldsymbol{\mu}_2$.

多元检验考虑了每组的两个测量值之间的正相关性, 而这一重要信息在一元检验中却不幸被忽略了. 上述 $T^2$ 检验等价于\hyperref[eq:6-42]{式(6-42)}的 MANOVA 检验. \hfill$\blacksquare$

\noindent\textbf{例 6.17}\quad (需要二元检验估计均值差的蜥蜴数据)

一个动物学家收集了美国西南部的蜥蜴. 他测量了蜥蜴的重量(克)和它们的鼻子到排泄口的长度(mm). 由于尾巴常常掉落, 这个长度更具代表性. 采集于 1997 年和 1999 年的两类蜥蜴 C 和 S 的数据见\hyperref[tab:6.7]{表 6.7}. 注意到 C 类蜥蜴有 20 个数据, S 类蜥蜴有 40 个数据.

\begin{mvtable}
\phantomsection\label{tab:6.7}\mvtabcap{表 6.7\quad 两类蜥蜴数据}
\begin{tabular}{cc|cc|cc}
\multicolumn{2}{c}{C} & \multicolumn{2}{c}{S} & \multicolumn{2}{c}{S}\\\midrule
重量&SVL&重量&SVL&重量&SVL\\\midrule
7.513&74.0&13.911&77.0&14.666&80.0\\
5.032&69.5&5.236&62.0&4.790&62.0\\
5.867&72.0&37.331&108.0&5.020&61.5\\
11.088&80.0&41.781&115.0&5.220&62.0\\
2.419&56.0&31.995&106.0&5.690&64.0\\
13.610&94.0&3.962&56.0&6.763&63.0\\
18.247&95.5&4.367&60.5&9.977&71.0\\
16.832&99.5&3.048&52.0&8.831&69.5\\
15.910&97.0&4.838&60.0&9.493&67.5\\
17.035&90.5&6.525&64.0&7.811&66.0\\
16.526&91.0&22.610&96.0&6.685&64.5\\
4.530&67.0&13.342&79.5&11.980&79.0\\
7.230&75.0&4.109&55.5&16.520&84.0\\
5.200&69.5&12.369&75.0&13.630&81.0\\
13.450&91.5&7.120&64.5&13.700&82.5\\
14.080&91.0&21.077&87.5&10.350&74.0\\
14.665&90.0&42.989&109.0&7.900&68.5\\
6.092&73.0&27.201&96.0&9.103&70.0\\
5.264&69.5&38.901&111.0&13.216&77.5\\
16.902&94.0&19.747&84.5&9.787&70.0\\
\end{tabular}
\end{mvtable}

\vspace{2pt}
\noindent\rule{\linewidth}{0.4pt}

\noindent\small SVL=鼻子到排泄口的长度

\noindent\small 资料来源: 数据经 Kevin E. Bonine 许可使用.
